\subsection{Egg ensemble}

Results for the standard DSI-ESMDA are detailed in \cref{tab:egg-dsi-esmda},
while results for the proposed Gated DSI-ESMDA under a mild gating intensity are shown in \cref{tab:egg-gated-dsi-esmda}.

\Cref{tab:egg-dsi-esmda} shows that DSI-ESMDA posteriors had a higher posterior coverage during the history period when localization was applied, 
in exchange for a slightly higher data mismatch average and standard deviation.
Indeed, averaging across groundtruths, 
localization increased the posterior coverage from 
35.4\% ($\pm$ 0.84\%) to 52.7\% ($\pm$ 2.11\%),
while the average data mismatch increased from 
0.496 ($\pm$ 0.005) to 0.555 ($\pm$ 0.02).
It is also clear that localization tends to increase posterior coverage in the validation period,
but it may also decrease it, as seen in the groundtruth 30 case.
%
While localization did not decrease posterior MSE,
values being close to 1 indicate that posterior can be considered good (MSE$\leq1$) or at least fair ($1<$MSE$\leq 2$)
\cite[section 9.4.2.2]{Emerick:2025}.

It is known \cite[appendix]{Oliver:2018} that the expected mismatch is 0.5
for randomized maximum likelihood samples.
Since the localization did not move the average mismatch significantly away from 0.5,
it is reasonable to say that the increase in posterior coverage is overall beneficial. 

From \cref{tab:egg-gated-dsi-esmda}, it can be seen that the proposed Gated DSI-ESMDA method, across groundtruths,
has an average posterior coverage in the history period of 41.00\% ($\pm$ 1.76\%) 
and an average mismatch of 0.513 ($\pm$ 0.008) when localization isn't applied,
versus 57.3\% ($\pm$ 1.25\%) and 0.598 ($\pm$ 0.01) when it is.
%
The increase in history coverages relative to those obtained by the DSI-ESMDA under the same conditions is significant,
although so is the increase in the data mismatch.
In the validation period, it is clear that localization has a positive effect on the posterior coverage,
and that the proposed method posterior coverages tend to be slightly higher than those obtained by the DSI-ESMDA.
%
Localization did not tend to decrease the posterior MSE, which remained good or fair. 
This was expected, since the same behavior was observed in the standard DSI-ESMDA.

The following short-hands are used in the tables and figures of this section:
Loc.: Localization is applied;
Pr.: Prior ensemble;
Po.: Posterior ensemble;
Avg.: Average value;
Std. Dev.: Standard deviation.

\begin{table}[htbp]
    \centering
    \IBGEtab{%
        \caption{%
            DSI-ESMDA results for the Egg ensemble.
        }
        \label{tab:egg-dsi-esmda}
    }{%
        \def\arraystretch{0.97}
        \begin{tabular}{ccccccc}
\toprule
\multirow{2}{*}{
    Gt.
}
&
\multirow{2}{*}{Loc.} 
& 
\multicolumn{2}{c}{OC} & 
\multirow{2}{*}{
    MSE
} & 
\multicolumn{2}{c}{Data mismatch} 
\\
& & His. & Val. &  & Avg. & Std. dev.
\\
\midrule
\multirow{2}{*}{10} & No & 0.34 & 1.00 & 0.96 & 0.50 & 5.1 $\times 10^{-3}$ \\
 & Yes & 0.48 & 1.00 & 0.96 & 0.55 & 0.03 \\
\midrule
\multirow{2}{*}{20} & No & 0.36 & 0.98 & 0.93 & 0.49 & 4.8 $\times 10^{-3}$ \\
 & Yes & 0.52 & 1.00 & 0.93 & 0.53 & 0.02 \\
\midrule
\multirow{2}{*}{30} & No & 0.36 & 1.00 & 0.94 & 0.49 & 4.6 $\times 10^{-3}$ \\
 & Yes & 0.53 & 0.90 & 0.94 & 0.57 & 0.03 \\
\midrule
\multirow{2}{*}{40} & No & 0.35 & 1.00 & 0.96 & 0.50 & 5.0 $\times 10^{-3}$ \\
 & Yes & 0.52 & 1.00 & 0.96 & 0.55 & 0.02 \\
\midrule
\multirow{2}{*}{50} & No & 0.35 & 0.90 & 0.94 & 0.50 & 5.2 $\times 10^{-3}$ \\
 & Yes & 0.53 & 1.00 & 0.94 & 0.55 & 0.03 \\
\midrule
\multirow{2}{*}{60} & No & 0.36 & 0.94 & 0.94 & 0.49 & 4.7 $\times 10^{-3}$ \\
 & Yes & 0.53 & 0.98 & 0.94 & 0.55 & 0.03 \\
\midrule
\multirow{2}{*}{70} & No & 0.34 & 1.00 & 0.95 & 0.50 & 5.6 $\times 10^{-3}$ \\
 & Yes & 0.53 & 1.00 & 0.95 & 0.55 & 0.02 \\
\midrule
\multirow{2}{*}{80} & No & 0.36 & 0.90 & 0.94 & 0.50 & 5.3 $\times 10^{-3}$ \\
 & Yes & 0.56 & 1.00 & 0.94 & 0.58 & 0.04 \\
\midrule
\multirow{2}{*}{90} & No & 0.36 & 0.96 & 0.95 & 0.50 & 4.9 $\times 10^{-3}$ \\
 & Yes & 0.55 & 1.00 & 0.95 & 0.59 & 0.05 \\
\midrule
\multirow{2}{*}{100} & No & 0.35 & 1.00 & 0.93 & 0.49 & 4.7 $\times 10^{-3}$ \\
 & Yes & 0.52 & 1.00 & 0.93 & 0.53 & 0.02 \\
\bottomrule
\end{tabular}

    }{%
        \fonte{Self-elaboration.}
    }
\end{table}
%
\begin{table}[htbp]
    \centering
    \IBGEtab{%
        \caption{%
            Gated DSI-ESMDA ($\lambda=2$) results for the Egg ensemble.
        }
        \label{tab:egg-gated-dsi-esmda}
    }{%
        \def\arraystretch{0.97}
        \begin{tabular}{ccccccc}
\toprule
\multirow{2}{*}{
    Gt.
}
&
\multirow{2}{*}{Loc.} 
& 
\multicolumn{2}{c}{OC} & 
\multirow{2}{*}{
    MSE
} & 
\multicolumn{2}{c}{Data mismatch} 
\\
& & His. & Val. &  & Avg. & Std. Dev.
\\
\midrule
\multirow{2}{*}{10} & No & 0.41 & 1.00 & 0.97 & 0.52 & 0.01 \\
 & Yes & 0.57 & 1.00 & 0.97 & 0.60 & 0.02 \\
\midrule
\multirow{2}{*}{20} & No & 0.43 & 1.00 & 0.94 & 0.51 & 0.01 \\
 & Yes & 0.58 & 1.00 & 0.94 & 0.58 & 0.02 \\
\midrule
\multirow{2}{*}{30} & No & 0.39 & 0.99 & 0.95 & 0.51 & 0.01 \\
 & Yes & 0.55 & 0.90 & 0.95 & 0.61 & 0.03 \\
\midrule
\multirow{2}{*}{40} & No & 0.42 & 1.00 & 0.97 & 0.53 & 0.01 \\
 & Yes & 0.57 & 1.00 & 0.97 & 0.60 & 0.02 \\
\midrule
\multirow{2}{*}{50} & No & 0.43 & 0.90 & 0.95 & 0.51 & 0.01 \\
 & Yes & 0.59 & 1.00 & 0.95 & 0.60 & 0.02 \\
\midrule
\multirow{2}{*}{60} & No & 0.41 & 0.95 & 0.94 & 0.51 & 0.01 \\
 & Yes & 0.57 & 0.96 & 0.95 & 0.60 & 0.02 \\
\midrule
\multirow{2}{*}{70} & No & 0.40 & 1.00 & 0.96 & 0.52 & 0.01 \\
 & Yes & 0.57 & 1.00 & 0.97 & 0.60 & 0.02 \\
\midrule
\multirow{2}{*}{80} & No & 0.40 & 0.93 & 0.95 & 0.50 & 0.01 \\
 & Yes & 0.59 & 1.00 & 0.95 & 0.60 & 0.03 \\
\midrule
\multirow{2}{*}{90} & No & 0.38 & 0.94 & 0.96 & 0.51 & 9.3 $\times 10^{-3}$ \\
 & Yes & 0.56 & 0.99 & 0.97 & 0.61 & 0.04 \\
\midrule
\multirow{2}{*}{100} & No & 0.43 & 1.00 & 0.94 & 0.51 & 9.3 $\times 10^{-3}$ \\
 & Yes & 0.58 & 1.00 & 0.94 & 0.58 & 0.02 \\
\bottomrule
\end{tabular}

    }{%
        \fonte{Self-elaboration.}
    }
\end{table}

\Cref{fig:egg-summaries} summarizes results for the Egg ensemble by averaging posterior coverage and MSE across groundtruths 
for different gating intensities, with null gating corresponding to standard DSI-ESMDA.

Gating seems to produce effects similar to those of localization, especially when its intensity is increased.
In the history period, there is a clear trend of increasing posterior coverage with increasing gating intensity,
at the cost of increased MSE. 
Importantly, the was little effect on validation coverage, but it did tend to increase relative to the prior regardless of gating intensity. 
This is likely an effect of the prior ensemble being a good representative of the groundtruths, 
as the FOPT CDF suggested.

\begin{figure}[htbp]
    % \centering
    \caption{%
        Average posterior coverage and MSE values found in the Egg ensemble.
    }
    \label{fig:egg-summaries}
    \begin{subfigure}{.49\linewidth}
        \centering
        \includegraphics[width=\linewidth]{results/egg/summary_plots/summary_history_coverage.pdf}
        \caption{History coverage.}
    \end{subfigure}
    %
    \begin{subfigure}{.49\linewidth}
        \centering
        \includegraphics[width=\linewidth]{results/egg/summary_plots/summary_validation_coverage.pdf}
        \caption{Validation coverage.}
    \end{subfigure}
    %
    \begin{subfigure}{.49\linewidth}
        \centering
        \includegraphics[width=\linewidth]{results/egg/summary_plots/summary_history_mse.pdf}
        \caption{MSE.}
    \end{subfigure}
    \fonte{Self-elaboration.}
\end{figure}